\documentclass[11pt]{article}

% Packages for formatting
\usepackage[utf8]{inputenc}
\usepackage[T1]{fontenc}
\usepackage{lmodern}
\usepackage{geometry}
\usepackage{graphicx}
\usepackage{amsmath}
\usepackage{amssymb}
\usepackage{caption}
\usepackage{color}
\usepackage{listings}
\usepackage{hyperref}
\usepackage{fancyhdr}
\usepackage{tcolorbox}
\usepackage{enumitem}
\usepackage{float}
\usepackage[table]{xcolor}

\geometry{margin=1in}
\pagestyle{fancy}
\fancyhf{}
\rhead{Deep Learning Experiments}

\rfoot{Page \thepage}

% Colors for code
\definecolor{codebg}{rgb}{0.95,0.95,0.95}
\definecolor{codegray}{rgb}{0.4,0.4,0.4}
\definecolor{codepurple}{rgb}{0.58,0,0.82}
\definecolor{codeblue}{rgb}{0.0,0.0,0.6}

% Python code style
\lstdefinestyle{python}{
    backgroundcolor=\color{codebg},   
    commentstyle=\color{codegray},
    keywordstyle=\color{codeblue}\bfseries,
    numberstyle=\tiny\color{gray},
    stringstyle=\color{codepurple},
    basicstyle=\ttfamily\footnotesize,
    breaklines=true,
    captionpos=b,
    keepspaces=true,
    numbers=left,
    numbersep=5pt,
    showspaces=false,
    showstringspaces=false,
    showtabs=false,
    tabsize=2,
    language=Python
}


\begin{document}


\section*{Experiment 1: 1D CNN}
\subsection*{About}
In this experiment, we implement a 1D Convolutional Neural Network (1D CNN) using PyTorch for classification of time-series data. Unlike 2D CNNs that operate over images, a 1D CNN processes sequential data where convolutional filters slide along one dimension, capturing local temporal patterns in signals such as audio, sensor data, or financial time series.

Mathematically, for an input sequence $x \in \mathbb{R}^{L}$ of length $L$, a 1D convolution operation with kernel $w \in \mathbb{R}^{k}$ of size $k$ and bias $b$ produces an output feature map $y \in \mathbb{R}^{L - k + 1}$ as:
\[
y[i] = \sum_{j=0}^{k-1} x[i + j] \cdot w[j] + b
\]
This operation is repeated with multiple filters to learn diverse features. Stacking multiple convolutional and pooling layers allows the model to capture hierarchical patterns over longer temporal contexts.

The 1D CNN model is trained to minimize a classification loss (e.g., cross-entropy) using stochastic gradient descent. We evaluate the model’s performance based on classification accuracy and loss convergence over training epochs.


\subsection*{Code}
\lstinputlisting[style=python, caption={1DCNN}]{1dcnn.py}%replace with your actual py file

\subsection*{Output}
\begin{tcolorbox}[colback=codebg, colframe=gray, title=1D CNN Training Output]

The 1D Convolutional Neural Network (CNN) was trained for 10 epochs on a sequential dataset for binary classification. The performance at the final epoch is summarized below:

\begin{itemize}
    \item \textbf{Epoch 10:} Train Accuracy = 99.25\%, Validation Accuracy = 89.50\%
\end{itemize}
\begin{figure}[H]
\centering
\includegraphics[width=0.7\linewidth]{1dcnn.png} 
% replace with your actual png file
\caption{Training and Validation Accuracy Curve for 1D CNN Model}
\end{figure}

\end{tcolorbox}

\subsection*{Observation}
\begin{itemize}
    \item The 1D CNN model achieved a very high training accuracy, indicating that it learned the training patterns effectively.
  
\end{itemize}
\newpage

\end{document}